% =============================================================================
%  4. Experimental setup
%  Data, metrics, protocols, reproducibility. Deliberately separate from
%  Section 3: the METHOD is what we built, the SETUP is what we measured it
%  with, and a reviewer checking whether a number is fair reads this section.
% =============================================================================
\section{Experimental Setup}
\label{sec:setup}

\subsection{Data}
\label{sec:setup:data}

The private evaluation set comprises 68 labelled cases over three trams and
nine camera views: 31 anomalous and 37 clean frames carrying 73 ground-truth
instances. \tabref{tab:data} gives the composition. Every case names the
reference it is compared against, and instance boxes are expressed in that
reference's pixel space.

\begin{table}[t]
\centering
\caption{Composition of the private evaluation set: 68 cases, 73 instances,
3 trams, 9 camera views.}
\label{tab:data}
\begin{tabular}{@{}llr@{\hspace{2em}}llr@{}}
\toprule
\multicolumn{3}{@{}l}{\textbf{Instances by type}} &
\multicolumn{3}{l@{}}{\textbf{Cases by provenance}} \\
\midrule
& forgotten object & 53 & & real camera frame        & 54 \\
& graffiti         &  7 & & inpainted anomaly        & 11 \\
& litter           &  7 & & self-staged              &  2 \\
& damage           &  6 & & viewpoint variant        &  1 \\
\cmidrule(r){2-3}\cmidrule(l){5-6}
& \textbf{total}   & \textbf{73} & & \textbf{total}  & \textbf{68} \\
\bottomrule
\end{tabular}
\end{table}

Eleven cases carry an anomaly composited into a real frame of the same camera
by image inpainting. They are labelled as such and exist because two of the
four anomaly types --- damage and graffiti --- are rare enough in a working
tram that waiting for them was not possible within the study period. Synthetic
augmentation of this kind is what Caetano et al.~\cite{Caetano_2022} recommend
for exactly this domain, where they report the lack of public in-vehicle data
as the binding constraint. People are
never instances: a kept box on a person counts as a false-positive region.

We state the obvious limitation once, here, and return to it in
\secref{sec:discussion}: this is a small, single-deployment, non-public
dataset. \secref{sec:results:cdnet} is the mitigation.

\subsection{Metrics}
\label{sec:setup:metrics}

Three levels are reported, because the central claim of the paper is precisely
that they move independently.

\paragraph{Frame level.} Each case is a binary decision --- a frame is flagged
when at least one region survives the \judge{}. We report $F_1$, recall and
specificity over the 68 cases.

\paragraph{Object level.} A ground-truth instance counts as detected if any
kept region overlaps it under a lenient rule ($\mathrm{IoU} > 0.1$ or centre
containment). A strict count at $\mathrm{IoU} \ge 0.3$ is reported alongside,
and \emph{the gap between the two is the box-quality signal} --- a proposal
covering most of the frame satisfies the lenient rule trivially and is useless
to a judge that only sees the crop. Kept regions matching no instance are
false-positive regions, and region precision is the fraction of kept regions
that match one.

\paragraph{Image-level AUROC.} For the comparisons against industrial
anomaly-detection models (\secref{sec:results:baselines},
\secref{sec:results:transfer}), which produce score maps rather than regions,
we use the standard metric of that literature,
%
\begin{equation}
  \mathrm{score}(x) = \operatorname{mean}\bigl( \text{top } 1\,\% \text{ of patch scores of } x \bigr),
  \qquad
  \text{I-AUROC} = \Pr\bigl[\, \mathrm{score}(x^{-}) > \mathrm{score}(x^{+}) \,\bigr],
  \label{eq:iauroc}
\end{equation}
%
with $x^{-}$ anomalous, $x^{+}$ clean, and masked patches excluded from the
top-$1\,\%$ pool.

\paragraph{One reporting rule.} Any recall is read next to the \judge{}'s YES
rate. With 73 instances among 654 proposed regions, a calibrated judge answers
YES about 11\,\% of the time; a judge answering YES to 88\,\% of crops also
recovers most instances, and reporting that as detection would be an artefact
of the proposal distribution rather than a property of the model. This follows
the sampling concern made explicit by POPE~\cite{Li_2023}.

\subsection{Per-camera sample sizes}
\label{sec:setup:sizes}

The AUROC comparisons are per camera, and their sample sizes differ by an
order of magnitude, so we state them once and let them bound every claim that
follows. Camera \texttt{tram\_1762} contributes 16 anomalous and 11 clean
frames (176 pairs); the four \texttt{3333} cameras contribute 14 anomalous and
7 clean frames in total (98 pairs), of which no individual camera exceeds 8
pairs; the three \texttt{1760} cameras have no labelled anomalies and
contribute score levels rather than AUROCs. Only the \texttt{tram\_1762}
column and the \texttt{3333}-pooled column carry enough pairs to support an
AUROC claim, and individual \texttt{3333} columns are reported for completeness
only.

\subsection{Public-data protocol}
\label{sec:setup:cdnet}

To evaluate the proposal stage on data we did not build, we run it on CDnet
2014~\cite{Wang_2014} across seven categories. The reference is the per-pixel
median of up to 50 frames preceding each sequence's \texttt{temporalROI},
which reproduces our clean reference by construction; every $k$-th ROI frame is
then evaluated. The official label semantics (static, hard shadow, outside ROI,
unknown, motion) and the spatial ROI are respected, and the ROI is resampled
with nearest-neighbour interpolation where a sequence ships an ROI of different
size from its frames.

The person veto is \textbf{disabled} for these runs. On CDnet the annotated
foreground \emph{is} people, so leaving it enabled would measure the veto
rather than the \proposer{}. Two views of each configuration are reported: the
per-pixel change mask, and the rasterised proposal boxes --- the latter being
what actually reaches the \judge{} --- both scored with the official CDnet
metrics.

We also report a qualitative check on ABODA, the one public dataset whose
anomaly class is ours. It ships eleven CCTV sequences and no annotations, so it
yields figures and not numbers; the event-level protocol of Luna et
al.~\cite{Luna_2018} would require the temporal annotation package we were
unable to obtain (\secref{sec:discussion}).

\subsection{Implementation and reproducibility}
\label{sec:setup:repro}

All experiments run on a single NVIDIA RTX~3080~Ti (12~GB). Models: DINOv2
ViT-S/14 with registers (frozen)~\cite{Oquab_2023}, YOLOv8-nano for the person
veto~\cite{Jocher_2023_yolov8}, a 9B GLM-V--family
judge~\cite{VTeam_2025_glm45v} in 4-bit via Ollama, and our Dinomaly
re-implementation~\cite{Guo_2025}. The judge sweep additionally covers
Qwen2.5-VL-7B~\cite{Bai_2025_qwen25vl}, Qwen3-VL-8B and
InternVL3.5-8B~\cite{Wang_2025_internvl35}. Baselines use the anomalib~2.6
reference implementations~\cite{Akcay_2022}.

Per-region judge verdicts are cached on a key covering image, reference, box,
model, prompt and occlusion mask, so a re-run reproduces a previous run exactly
while a changed mask or prompt correctly invalidates it. The cache key
deliberately excludes the \proposer{}: a box is a box, so identical coordinates
from two \proposer{}s share one verdict --- which is also what makes the sweep
of \secref{sec:results:main} affordable.

Code, ground truth, per-run reports and a decision log are available in the
project repository. The tram footage itself cannot be released, which is the
reason \secref{sec:results:cdnet} exists.
